% !TeX root = ../../Book1.tex
\chapter{测度的微分与 Lebesgue 点}
\label{b1:ch:14}
% 本章摘要（不计入编号）
\begin{chapterabstract}
Radon--Nikodym 定理用积分等式给出密度，却没有直接说明怎样在一个点附近计算它。本章把空间中的小球引入这一问题：先比较球内质量与体积，再把分母换成任意 Radon 测度。覆盖定理与极大不等式保证这些局部比值几乎处处收敛，并解释为何它们恢复的是绝对连续部分，而不能自动找回整个奇异测度。

对局部可积函数，微分定理进一步给出 Lebesgue 点：不仅平均值恢复函数值，平均绝对偏差也趋于零。围绕这一性质，我们讨论积分等价类的精确代表、改变观察集合或积分核后的恢复结果，以及只控制坏点比例的近似连续性。通常平均、强平均值与近似极限之间的区别，是本章需要始终保留的一条界线。
\end{chapterabstract}

% 本章目标（不计入编号）
\begin{learninggoals}
读完本章，应能核对球比值的定义域与可测性，说明开球、闭球和几乎处处的基准测度；能用极大估计把连续函数的结论延拓到局部可积函数，并在任意 Radon 测度下完成绝对连续与奇异两部分的微分论证；能区分特定代表的点值与由积分确定的值，判断什么时候可以更换局部平均的形状或核；最后，能通过例子说明近似连续性为何弱于 Lebesgue 点性质。本章只预告近似微分，不提前假定一般可测函数具有几乎处处的线性近似。
\end{learninggoals}

% !TeX root = ../../Book1.tex
\section{密度与测度比}
\label{b1:sec:1401}

导数用越来越小的增量描述局部变化。对测度而言，可以把增量换成围绕一点缩小的球，再比较球中的质量与体积。若测度原本由函数积分得到，这个比值应当恢复函数的局部值；若质量集中在体积为零的集合上，同一个比值则可能趋于无穷，或者在几乎所有点都看不见这部分质量。

本节先定义这些比值，说明它们何时有意义，以及哪些结论只需测度连续性便能得到。几乎处处存在极限是后面两节的任务。集合密度的初步例子可对照 Axler 的《Measure, Integration \& Real Analysis》\cite[4.22--4.24]{Axler2020Measure}；相对于一般 Radon 测度的球平均及支集条件，可结合 Fremlin 的《Measure Theory》\cite[472D--E]{Fremlin2016Measure}阅读。后者用闭球，本章的测度比也采用这一约定。

\subsection{上、下密度}
\label{b1:subsec:140101}

本节固定在 \(\mathbb R^d\) 上讨论，\(m\) 表示 Lebesgue 测度。由伸缩性质，正半径球的体积有限且为正。对 Radon 测度 \(\nu\)，有界闭球的质量也有限，因此下面每个固定半径的比值都是非负实数；趋于零尺度时的极限却可以为无穷。

% 实读来源：Axler2020Measure，4.22--4.24，印刷 pp.112--113；
% Simon1983Lectures，§3 开头及 Remark 3.1，印刷 p.10。
% 改编：本节先采用环境维数的 Lebesgue 体积归一化，不提前引入 Hausdorff 测度；集合例子与振荡环带计算独立展开。
\begin{definition}
\label{b1:def:measure-densities}
设 \(\nu\) 是 \(\mathbb R^d\) 上的正 Radon 测度。对每个 \(x\in\mathbb R^d\)，把
\[
\underline D_m\nu(x)
=\liminf_{r\downarrow0}
\frac{\nu\bigl(\overline B(x,r)\bigr)}{m\bigl(\overline B(x,r)\bigr)},
\quad
\overline D_m\nu(x)
=\limsup_{r\downarrow0}
\frac{\nu\bigl(\overline B(x,r)\bigr)}{m\bigl(\overline B(x,r)\bigr)}
\]
分别称为 \(\nu\) 在 \(x\) 处相对于 Lebesgue 测度的\term[xia-midu]{下密度}[lower density]与\term[shang-midu]{上密度}[upper density]。若二者相等，则把共同值记作 \(D_m\nu(x)\)，称为密度，允许其值为 \(+\infty\)。
\end{definition}

上、下密度记录的并非沿某一条预选半径序列的极限。例如，对比值函数 \(a(r)\)，有
\[
\limsup_{r\downarrow0}a(r)
=\inf_{\delta>0}\sup_{0<r<\delta}a(r),
\quad
\liminf_{r\downarrow0}a(r)
=\sup_{\delta>0}\inf_{0<r<\delta}a(r).
\]
只计算 \(r=2^{-n}\) 上的值，一般还不足以确定任意半径趋零时的行为。

\begin{definition}
\label{b1:def:density-point}
设 \(E\subseteq\mathbb R^d\) 为 Lebesgue 可测集。在上一定义中取 \(\nu=m\llcorner E\)，得到
\[
\begin{aligned}
\underline\theta(E,x)
&=\liminf_{r\downarrow0}
\frac{m\bigl(E\cap\overline B(x,r)\bigr)}{m\bigl(\overline B(x,r)\bigr)},\\
\overline\theta(E,x)
&=\limsup_{r\downarrow0}
\frac{m\bigl(E\cap\overline B(x,r)\bigr)}{m\bigl(\overline B(x,r)\bigr)}.
\end{aligned}
\]
若两者相等，则把共同值记作 \(\theta(E,x)\)。若 \(\theta(E,x)=1\)，则称 \(x\) 是 \(E\) 的一个\term[midu-dian]{密度点}[density point]。
\symidx{\(\theta(E,x)\)：集合在一点的密度}
\end{definition}

集合的上、下密度都在 \([0,1]\) 内，并且补集交换上、下极限：
\[
\underline\theta(E^{\mathrm c},x)=1-\overline\theta(E,x),
\quad
\overline\theta(E^{\mathrm c},x)=1-\underline\theta(E,x).
\]
这些等式来自每个球中 \(E\) 与其补集的质量之和等于整个球的质量。若 \(m(E\mathbin{\triangle} F)=0\)，则 \(E\) 与 \(F\) 在每个球中的质量相同，所以它们在每一点的上、下密度都相同，而不只是几乎处处相同。

内点一定是密度点，因为足够小的闭球全在集合内；若 \(x\notin\overline E\)，则足够小的球与 \(E\) 不交，密度为零。反过来，密度点不必是内点，甚至不必属于原集合。例如 \(\mathbb R^d\setminus\{0\}\) 在原点的密度仍为 \(1\)。集合 \(\mathbb Q^d\) 的闭包是整个空间，但它在每一点的密度都为零。这些例子说明密度忽略零测修改，而通常的内点与闭包会被零测修改改变。

边界点也可能具有严格介于 \(0\) 与 \(1\) 之间的密度。取半空间 \(H=\{\,x\in\mathbb R^d\mid x_d>0\,\}\)。在边界上一点 \(x\) 处，关于超平面的反射把球的两半互换；由\cref{b1:lem:coordinate-hyperplane-null}及反射不变性，分界超平面为零测集，两半球体积相等。因此 \(\theta(H,x)=1/2\)。

还可以让密度完全没有极限。令 \(r_n=2^{-2^n}\)，并取
\[
E=\bigcup_{k=0}^{\infty}
\{\,x\in\mathbb R^d\mid r_{2k+1}<|x|<r_{2k}\,\}.
\]
这是一串交替保留与删去的同心环带。记 \(a(r)\) 为 \(E\) 在原点半径 \(r\) 球中的体积比例。球面的零测性与伸缩公式给出
\[
a(r_{2k})\geq1-\left(\frac{r_{2k+1}}{r_{2k}}\right)^d,
\quad
a(r_{2k+1})\leq\left(\frac{r_{2k+2}}{r_{2k+1}}\right)^d.
\]
由于 \(r_{n+1}/r_n=2^{-2^n}\longrightarrow0\)，第一列比例趋于 \(1\)，第二列趋于 \(0\)。结合 \(0\leq a(r)\leq1\)，可知 \(\overline\theta(E,0)=1\)、\(\underline\theta(E,0)=0\)。

这些都是逐点计算，尚未限制异常点集的大小。\cref{b1:thm:lebesgue-density}将证明：对任意 Lebesgue 可测集 \(E\)，几乎每个属于 \(E\) 的点都是密度点，而在补集内密度几乎处处为零。该结论将在下一节从积分的微分定理推出，这里不把它作为已知性质使用。中文伴读可参看那汤松的《实变函数论》第5版中译本\cite[第九章第6节]{Natanson2010RealChinese}，该书以“稠密点、近似连续”为这一话题的节名；本书使用“密度点”，将它与拓扑意义下的稠密性分开。

\subsection{密度比}
\label{b1:subsec:140102}

若测度集中在曲线或离散点上，Lebesgue 体积未必是合适的基准。改用另一个测度 \(\mu\) 后，首先要处理分母可能为零的问题。

\begin{definition}
\label{b1:def:measure-support}
设 \(\mu\) 为 \(\mathbb R^d\) 上的正 Radon 测度。集合
\[
\operatorname{supp}\mu
=\{\,x\in\mathbb R^d\mid \mu\bigl(B(x,r)\bigr)>0\ \forall r>0\,\}
\]
称为\term[cedu-zhiji]{测度的支集}[support of a measure]。
\symidx{\(\operatorname{supp}\mu\)：测度的支集}
\end{definition}

支集是闭集，其补集为零测集。事实上，若 \(\mu\bigl(B(x,r)\bigr)=0\)，则 \(B(x,r/2)\) 中每一点都有一个零测开球邻域，因而都不在支集中，所以支集的补集开。再取欧氏空间的一个可数开球基。每个不在支集中的点都落在某个包含于零测开球的基元素内；所有这样的基元素组成可数族，测度均为零，并覆盖支集的补集。可数次可加性便给出
\[
\mu\bigl(\mathbb R^d\setminus\operatorname{supp}\mu\bigr)=0.
\]
这里使用了欧氏空间的第二可数性，不能仅凭不可数个零测开集的并就作出零测结论。

给定两个正 Radon 测度 \(\mu,\nu\)，对 \(x\in\operatorname{supp}\mu\) 定义球上的质量比
\[
Q_r^{\mu,\nu}(x)
=\frac{\nu\bigl(\overline B(x,r)\bigr)}{\mu\bigl(\overline B(x,r)\bigr)},
\quad r>0.
\]
由于闭球包含同中心开球，分母严格为正；由于闭球紧，分子与分母都有限。在支集外不指定这个比值，避免把 \(0/0\) 当成某个固定值来掩盖定义域。测度 \(\mu=0\) 的情形下支集为空，与此约定一致。

\begin{proposition}
\label{b1:prop:open-closed-density-ratio}
固定 \(x\in\operatorname{supp}\mu\)。在质量比中把闭球换成开球，不改变 \(r\downarrow0\) 时的上极限或下极限。这不要求任何球面为 \(\mu\)-零集或 \(\nu\)-零集。
\end{proposition}
\begin{proof}
对每个固定的 \(r>0\)，从 \(s>r\) 递减到 \(r\) 时，开球 \(B(x,s)\) 递减到 \(\overline B(x,r)\)。两测度在一个较大的闭球上都有限，因此测度从上连续性给出
\[
\frac{\nu\bigl(B(x,s)\bigr)}{\mu\bigl(B(x,s)\bigr)}
\longrightarrow Q_r^{\mu,\nu}(x).
\]
从 \(t<r\) 递增到 \(r\) 时，闭球 \(\overline B(x,t)\) 递增到 \(B(x,r)\)，同样有
\[
Q_t^{\mu,\nu}(x)
\longrightarrow
\frac{\nu\bigl(B(x,r)\bigr)}{\mu\bigl(B(x,r)\bigr)}.
\]
两次取商的极限都有严格正的极限分母。给定 \(\delta>0\)，当 \(r<\delta\) 时，上述逼近半径也可全取小于 \(\delta\)。所以两类比值在 \(0<r<\delta\) 上具有相同的上确界与下确界。再令 \(\delta\downarrow0\) 即得结论。
\end{proof}

固定半径的开球与闭球可能有不同质量；命题比较的是全部小半径组成的极限过程。对于原子恰好落在球面上的测度，这个区别不能略去。

\subsection{测度相对于测度的导数}
\label{b1:subsec:140103}

\begin{definition}
\label{b1:def:relative-density}
设 \(\mu,\nu\) 为正 Radon 测度。对 \(x\in\operatorname{supp}\mu\)，定义
\[
\underline D_\mu\nu(x)=\liminf_{r\downarrow0}Q_r^{\mu,\nu}(x),
\quad
\overline D_\mu\nu(x)=\limsup_{r\downarrow0}Q_r^{\mu,\nu}(x).
\]
它们分别是 \(\nu\) 相对于 \(\mu\) 的下密度与上密度。若二者相等，则把共同值记作 \(D_\mu\nu(x)\)，称为 \(\nu\) 相对于 \(\mu\) 在 \(x\) 处的\term[cedu-daoshu]{测度导数}[derivative of a measure]，允许取值 \(+\infty\)。
\symidx{\(\underline D_\mu\nu\)：测度比的下密度}
\symidx{\(\overline D_\mu\nu\)：测度比的上密度}
\symidx{\(D_\mu\nu\)：测度相对于测度的导数}
\end{definition}

当 \(\mu=m\) 时，这就回到第一小节的密度。这里的导数由空间中的小球定义，暂时还不是第\ref{b1:sec:0604}节中由积分等式确定的 Radon--Nikodym 导数。将二者联系起来，既要证明球比值几乎处处收敛，也要辨别奇异部分在极限中的表现。

% 实读来源：Fremlin2016Measure，卷4，472D--E 及习题 472Xd、472Yg，
% tmp/ch13-fremlin47.tex。只用作支集、可测性问题和测度比值识别的来源定位；
% 下列有理半径约化、开闭球比较、连续点与原子计算均在正文独立证明，不提前调用 472D 的结论。
\begin{proposition}
\label{b1:prop:relative-density-borel}
在 \(\operatorname{supp}\mu\) 上，固定半径的质量比 \(Q_r^{\mu,\nu}\) 是 Borel 可测函数。上、下密度也是扩展实值 Borel 可测函数；它们的定义中可把半径限制为正有理数。导数存在的点集以及导数存在且有限的点集均为 Borel 集，导数在这些定义域上 Borel 可测。
\end{proposition}
\begin{proof}
先看任意正 Radon 测度 \(\lambda\) 的球质量。若 \(x_n\to x\)，则对每个 \(\varepsilon>0\)，当 \(n\) 足够大时有
\(\overline B(x_n,r)\subseteq\overline B(x,r+\varepsilon)\)。于是
\[
\limsup_{n\to\infty}\lambda\bigl(\overline B(x_n,r)\bigr)
\leq\lambda\bigl(\overline B(x,r+\varepsilon)\bigr).
\]
让 \(\varepsilon\downarrow0\)，由较大闭球测度有限及从上连续性，右边趋于 \(\lambda\bigl(\overline B(x,r)\bigr)\)。因此固定半径的球质量是上半连续函数，特别是 Borel 可测的。取 \(\lambda=\mu,\nu\)，再在支集上对严格正的分母取商，得到 \(Q_r^{\mu,\nu}\) 的可测性。

固定 \(x\) 与 \(r>0\)，取有理数 \(q_j\downarrow r\)。两测度从上连续性给出
\(Q_{q_j}^{\mu,\nu}(x)\to Q_r^{\mu,\nu}(x)\)。若 \(r<1/n\)，可以令所有 \(q_j<1/n\)，所以在每个小半径区间上，有理半径给出相同的上确界与下确界。因此
\[
\begin{aligned}
\overline D_\mu\nu
&=\inf_{n\geq1}\sup_{\substack{q\in\mathbb Q\\0<q<1/n}}Q_q^{\mu,\nu},\\
\underline D_\mu\nu
&=\sup_{n\geq1}\inf_{\substack{q\in\mathbb Q\\0<q<1/n}}Q_q^{\mu,\nu}.
\end{aligned}
\]
右边都是可数次上下确界，故为 Borel 函数。导数存在的点集是两者相等之处：其补集可写为
\[
\bigcup_{a\in\mathbb Q}
\{\,x\in\operatorname{supp}\mu\mid
\underline D_\mu\nu(x)<a<\overline D_\mu\nu(x)\,\},
\]
因而可测；再与 \(\{\overline D_\mu\nu<\infty\}\) 相交便得到有限导数的定义域。在导数存在处，它等于上密度的限制，所以同样可测。
\end{proof}

最直接的恢复公式发生在密度函数连续的点上。

\begin{proposition}
\label{b1:prop:continuous-density-recovery}
设 \(\nu=f\mu\) 是正 Radon 测度，其中 \(f\) 非负可测。若 \(x\in\operatorname{supp}\mu\)，且所选代表 \(f\) 在 \(x\) 连续并取有限值，则
\[
D_\mu\nu(x)=f(x).
\]
\end{proposition}
\begin{proof}
给定 \(\varepsilon>0\)，由连续性可取 \(\delta>0\)，使 \(|f(y)-f(x)|<\varepsilon\) 对所有 \(|y-x|<\delta\) 成立。当 \(0<r<\delta\) 时，
\[
\begin{aligned}
\left|Q_r^{\mu,\nu}(x)-f(x)\right|
&\leq\frac{1}{\mu\bigl(\overline B(x,r)\bigr)}
\int_{\overline B(x,r)}|f(y)-f(x)|\,\mathrm d\mu(y)\\
&\leq\varepsilon.
\end{aligned}
\]
令 \(r\downarrow0\) 即得结论。
\end{proof}

这个论证不需要 \(\mu\) 与 Lebesgue 测度有什么关系，只需要球中分母为正，以及 \(f\) 在目标点的连续性。后面的微分定理要把连续性替换为可积性，并相应地把“该点成立”换成“几乎处处成立”。

\begin{proposition}
\label{b1:prop:atomic-density-ratio}
若 \(\mu(\{x\})>0\)，则
\[
D_\mu\nu(x)=\frac{\nu(\{x\})}{\mu(\{x\})}.
\]
\end{proposition}
\begin{proof}
当 \(r\downarrow0\) 时，闭球 \(\overline B(x,r)\) 递减到 \(\{x\}\)。两测度在一个固定闭球上有限，所以球的质量分别趋于 \(\nu(\{x\})\) 与 \(\mu(\{x\})\)。后者严格为正，故可以直接取商的极限。
\end{proof}

例如，取 \(\mu=\delta_0+m\)、\(\nu=2\delta_0+3m\)。在原点，测度导数为 \(2\)；在任意 \(x\ne0\) 处，足够小的球避开原点，质量比恒为 \(3\)。此例中的原点不是可任意修改的 \(\mu\)-零测点，因而密度函数在原点的值也由测度决定。

另一方面，若分母为 \(m\)、分子为 \(\delta_0\)，则原点处的比值为 \(1/m\bigl(\overline B(0,r)\bigr)\)，趋于 \(+\infty\)；在其他点处，小球质量最终为零。因此 \(D_m\delta_0\) 在 \(m\)-几乎处处等于零。把 \(\delta_0\) 换成 \(2\delta_0\)，甚至会得到处处相同的扩展实值导数，而两测度并不相等。这说明测度比的极限不能未经条件地代回积分，宣称它恢复整个原测度。\cref{b1:sec:1403}将用 Lebesgue 分解精确说明它恢复哪一部分。

% !TeX root = ../../Book1.tex
\section{Lebesgue 微分定理}
\label{b1:sec:1402}

一个函数的积分不受零测集上的改动影响，所以不可能由积分恢复任意指定的每一个点值。合理的问题是：局部平均能否在几乎处处恢复某个可测代表？对连续函数，这由连续性直接保证；对局部可积函数，需要将连续逼近的整体误差转换成平均的逐点误差。上一章的极大不等式正是这一步所需的控制。

本节采用连续稠密类与极大估计的路线，可对照 Axler 的《Measure, Integration \& Real Analysis》\cite[4.10]{Axler2020Measure}。先证明通常平均的微分，再把它加强为绝对偏差的平均趋零。两者不能只凭积分符号相似就视为同一陈述。

\subsection{局部平均}
\label{b1:subsec:140201}

设 \(\Omega\subseteq\mathbb R^d\) 开，\(f\in L^1_{\mathrm{loc}}(\Omega)\)。在 \(B(x,r)\) 的闭包包含于 \(\Omega\) 时，记
\[
 A_rf(x)=\frac1{\omega_dr^d}\int_{B(x,r)}f(y)\,\mathrm dy.
\]
\symidx{\(A_rf\)：中心球上的积分平均}
这里的 \(f\) 可以取实值或复值。球面为 Lebesgue 零测集，故平均与上一节的闭球密度约定相容。

对每个固定 \(r>0\)，\(A_rf\) 在允许的中心区域连续。证明与\cref{b1:prop:maximal-elementary}相同：中心平移时，球的示性函数在一个零测球面之外收敛，被积函数则由一个固定紧集上的 \(|f|\) 控制。对固定 \(x\)，\(r\mapsto A_rf(x)\) 也连续：若 \(r_n\to r>0\)，球的示性函数仍几乎处处收敛，分母也趋于一个正数。

因此，对 \(x\) 处所有小半径取上、下极限，可以只使用正有理半径。比如，选定一个处处有限的可测代表后，
\[
 \mathcal E f(x)=\limsup_{r\downarrow0}|A_rf(x)-f(x)|
\]
是可测函数。具体地，它可以写为对 \(n\) 取下确界、再对所有允许的 \(q\in\mathbb Q\cap(0,1/n)\) 取上确界；在不允许的中心处把相应的非负项延为零即可。每个内层都是可数上确界，且每个 \(x\in\Omega\) 都有充分小的允许半径。这一可测性说明避免了对不可数个半径分别选取例外零测集。

通常平均的收敛允许抵消。在实直线上，\(f(y)=\mathbf1_{(0,\infty)}(y)-\mathbf1_{(-\infty,0)}(y)\)，并取 \(f(0)=0\)，则 \(A_rf(0)=0\) 对所有 \(r\) 成立，但
\[
 \frac1{2r}\int_{-r}^r|f(y)-f(0)|\,\mathrm dy=1.
\]
所以“平均趋于点值”和“平均绝对偏差趋零”有不同的逐点含义。

\subsection{极大函数方法}
\label{b1:subsec:140202}

先在整个 \(\mathbb R^d\) 上设 \(f\in L^1\)。任取 \(g\in C_c(\mathbb R^d)\)，对每个半径都有
\[
 |A_rf(x)-f(x)|
 \leq A_r|f-g|(x)+|A_rg(x)-g(x)|+|g(x)-f(x)|.
\]
连续性使中间一项随 \(r\downarrow0\) 趋零，而第一项对所有 \(r\) 同时受 \(M(f-g)(x)\) 控制。因此
\begin{equation}\label{b1:eq:differentiation-maximal-error}
 \mathcal E f(x)\leq M(f-g)(x)+|f(x)-g(x)|.
\end{equation}
对 \(t>0\)，由 Hardy--Littlewood 弱型估计和 Chebyshev 不等式，得到
\begin{equation}\label{b1:eq:differentiation-exceptional-estimate}
 m_d\bigl(\{\mathcal E f>2t\}\bigr)
 \leq\frac{5^d+1}{t}\|f-g\|_1.
\end{equation}

这个估计没有要求 \(f-g\) 逐点一致小，也没有要求 \(M(f-g)\) 可积。它只要求大值位置的测度受 \(L^1\) 误差控制。左端已经包含任意小的全部尺度，并且不依赖所选的 \(g\)；这才允许用稠密性把右端送到零。

若对每个半径单独估计 \(A_r(f-g)\)，再说这些误差都很小，仍缺少统一半径的论证。另一方面，也不能直接对 \(r\downarrow0\) 套用控制收敛定理：归一化的球核越来越高，并没有一个显然可积的共同控制函数。极大估计处理的恰好是这两种直接论证留下的缺口。

\subsection{Lebesgue 微分定理}
\label{b1:subsec:140203}

\term*[lebesgue-weifen-dingli]{Lebesgue 微分定理}[Lebesgue differentiation theorem]说明，局部可积性已经足以保证通常平均在几乎处处恢复函数。其结论涉及固定可测代表，但任何两个代表得到的结论只会在一个零测集上不同。

% 实读来源：Axler2020Measure，4.10，印刷pp.108--109。
% 改编：全尺度有理半径约化、欧氏d维弱型控制、开集紧耗尽和复值版本均明确展开。
\begin{theorem}[Lebesgue 微分]\label{b1:thm:lebesgue-differentiation}
设 \(\Omega\subseteq\mathbb R^d\) 开，\(f\in L^1_{\mathrm{loc}}(\Omega)\)。则对几乎每个 \(x\in\Omega\)，
\[
 \lim_{r\downarrow0}\frac1{\omega_dr^d}\int_{B(x,r)}f(y)\,\mathrm dy=f(x).
\]
\end{theorem}
\begin{proof}
先设 \(\Omega=\mathbb R^d\)、\(f\in L^1\)。由\cref{b1:prop:lp-continuous-dense}，存在 \(g_j\in C_c\) 使 \(\|f-g_j\|_1\to0\)。将其代入\eqref{b1:eq:differentiation-exceptional-estimate}，可知对每个 \(t>0\)，\(m_d(\{\mathcal E f>2t\})=0\)。取 \(t=1/k\) 并作可数并，得到 \(\mathcal E f=0\) 几乎处处，即所求极限。

一般开集上，令
\[
 K_n=\{\,x\in\Omega:|x|\leq n,\quad
               \operatorname{dist}(x,\Omega^c)\geq1/n\,\}.
\]
约定到空集的距离为 \(\infty\)。\(K_n\) 紧，且 \(\Omega=\bigcup_nK_n\)。若 \(K_n\ne\varnothing\)，令 \(\rho_n=1/(3n)\)，并设
\[
 H_n=\{\,y:\operatorname{dist}(y,K_n)\leq\rho_n\,\}.
\]
它是包含于 \(\Omega\) 的紧集。把 \(f\mathbf1_{H_n}\) 延为整个空间上的零函数，记为 \(f_n\)，则 \(f_n\in L^1(\mathbb R^d)\)。对 \(x\in K_n\)、\(0<r<\rho_n\)，有 \(A_rf_n(x)=A_rf(x)\)，且 \(f_n(x)=f(x)\)。已证情形因而在 \(K_n\) 的一个零测集之外给出结论。最后对这些零测集作可数并，覆盖整个 \(\Omega\)。空的 \(K_n\) 不需处理。
\end{proof}

\begin{corollary}[Lebesgue 密度]\label{b1:thm:lebesgue-density}
若 \(E\subseteq\mathbb R^d\) 为 Lebesgue 可测集，则
\[
 \lim_{r\downarrow0}\frac{m_d\bigl(E\cap B(x,r)\bigr)}{\omega_dr^d}
 =\mathbf1_E(x)
\]
对几乎每个 \(x\in\mathbb R^d\) 成立。特别地，记 \(E^{(1)}\) 为 \(E\) 的所有密度点，则 \(m_d(E\mathbin{\triangle} E^{(1)})=0\)。
\end{corollary}
\begin{proof}
示性函数 \(\mathbf1_E\) 总是局部可积，无须假设 \(m_d(E)<\infty\)。将它代入微分定理即得第一式。由上一节密度函数的可测性，\(E^{(1)}\) 可测；第一式又说明 \(E\setminus E^{(1)}\) 与 \(E^{(1)}\setminus E\) 都是零测集。
\end{proof}

密度点无需是内点。一个正测度、内点为空的闭集仍在几乎每个自身点具有密度 \(1\)；第三章的正测度 Cantor 型集合便属此类。定理也没有说每个边界点密度均为 \(0\) 或 \(1\)：半空间的边界点密度为 \(1/2\)，而上一节的环带例子甚至没有密度极限。它限制的是这些异常行为能够占据的 Lebesgue 测度。

\subsection{\texorpdfstring{\(L^p_{\mathrm{loc}}\)}{Lp} 版本}
\label{b1:subsec:140204}

通常平均的极限还可以加强。关键是同时微分一个可数的辅助函数族，再利用值域中的稠密点逼近待恢复的数值。

\begin{theorem}[Lebesgue 微分，\texorpdfstring{\(L^p_{\mathrm{loc}}\)}{Lp loc} 偏差形式]\label{b1:thm:lp-lebesgue-differentiation}
设 \(1\leq p<\infty\)、\(f\in L^p_{\mathrm{loc}}(\Omega)\)。则对几乎每个 \(x\in\Omega\)，
\[
 \lim_{r\downarrow0}\frac1{\omega_dr^d}
          \int_{B(x,r)}|f(y)-f(x)|^p\,\mathrm dy=0.
\]
\end{theorem}
\begin{proof}
在标量域中取可数稠密集 \(D\)：实情形取 \(\mathbb Q\)，复情形取 \(\mathbb Q+i\mathbb Q\)。对每个 \(a\in D\)，函数 \(h_a=|f-a|^p\) 局部可积，因为
\[
 |f-a|^p\leq2^{p-1}\bigl(|f|^p+|a|^p\bigr).
\]
微分定理给出 \(A_rh_a(x)\to h_a(x)\) 几乎处处。删去关于可数个 \(a\) 的全部例外集，以及 \(f\) 不取有限值的零测集。固定余下的 \(x\)，在球上的归一化概率测度下应用 Minkowski 不等式，得到
\[
 \bigl(A_r|f-f(x)|^p(x)\bigr)^{1/p}
 \leq\bigl(A_r|f-a|^p(x)\bigr)^{1/p}+|a-f(x)|.
\]
对 \(r\downarrow0\) 取上极限，右端趋于 \(2|a-f(x)|\)。由 \(D\) 稠密，这个数可以任意小，故左端趋零。取 \(p\) 次幂得到结论。
\end{proof}

当 \(p=1\) 时，得到平均绝对偏差趋零，强于只知道 \(A_rf(x)\to f(x)\)。这也是下一节定义 Lebesgue 点的依据。当 \(p>1\) 时，Hölder 不等式说明 \(p\) 次偏差趋零蕴含一次偏差趋零，反向对某个指定点则未必成立。

还有一种不同的 \(L^p\) 陈述：对每个紧集 \(K\subseteq\Omega\)，
\begin{equation}\label{b1:eq:local-average-lp-convergence}
 \|A_rf-f\|_{L^p(K)}\longrightarrow0\quad(r\downarrow0).
\end{equation}
这同样对 \(1\leq p<\infty\) 成立。取 \(K\) 的一个仍在 \(\Omega\) 内的紧邻域 \(H\)，把 \(f\mathbf1_H\) 延为零。对充分小的 \(r\)，其球平均在 \(K\) 上与 \(A_rf\) 一致；而全空间的平均是核 \(\mathbf1_{B(0,1)}/\omega_d\) 生成的卷积逼近，故由\cref{b1:prop:approximate-identity-lp}得到\eqref{b1:eq:local-average-lp-convergence}。

点态偏差极限与局部范数收敛涉及不同的量词。前者固定中心再缩小球，后者在一整块中心区域上积分误差。对 \(p=1\) 的后者不能直接拿 \(Mf\) 作可积控制，因为上一章已看到极大算子并不在 \(L^1\) 上强型有界。对于 \(L^\infty_{\mathrm{loc}}\) 函数，一次乃至每个有限次幂的偏差仍几乎处处趋零，但不能据此断言局部本质上确界的偏差趋零。

% !TeX root = ../../Book1.tex
\section{Radon 测度的微分}
\label{b1:sec:1403}

% 实际读取：Fremlin2016Measure，作者官方 chap47.tex，472D--E 完整正文。
% 472D 的有理阈值、开集逼近和 Radon 细覆盖路线在本节完整改写；
% 原书的完备 Radon 约定在这里改为本书的 Borel 版本，坏集可测性另作说明。
% 472E 的一般 Radon 球比值估计只作方法对照，本节不把它作为未证前提。
% 奇异部分用13.2细覆盖直接证明；局部有限分解由第6章sigma有限RN定理构造。

Lebesgue 微分定理以体积作为分母。若质量集中在曲线、稀疏集合或原子上，
相对于体积的平均不能描述这种质量自身的局部行为。
本节同时允许分子和分母为任意 Radon 测度，证明小球质量之比仍能恢复绝对连续部分的密度。
决定这一结论的几何工具是上一章的 Radon 细覆盖，而不是球体积的缩放公式。

\subsection{Radon--Nikodym 导数的点态表示}
\label{b1:subsec:140301}

以下 \(\mu,\nu\) 都是 \(\mathbb R^d\) 上的正 Radon 测度，沿用第\ref{b1:sec:1102}节的 Borel 约定。
局部有限性保证每个有界闭球的质量有限；用 \(\overline B(0,n)\) 穷竭全空间，又知这两个测度都为 \(\sigma\)-有限。
由\cref{b1:prop:radon-outer-regularity}，它们在每个 Borel 集上也具有开外正则性。
这些性质不要求总质量有限。

记 \(S_\mu=\operatorname{supp}\mu\)。第\ref{b1:sec:1401}节已经说明
\(\mu(\mathbb R^d\setminus S_\mu)=0\)，并且在 \(x\in S_\mu\)、\(r>0\) 时有
\[
 0<\mu\bigl(\overline B(x,r)\bigr)<\infty.
\]
因此本节所有相对于 \(\mu\) 的微分极限，首先在 \(S_\mu\) 上取闭球比值；
其外无需为 \(0/0\) 规定数值，也不影响任何 \(\mu\)-几乎处处结论。
记号 \(D_\mu\nu\) 及其上、下极限沿用\cref{b1:def:relative-density}。
本节不假设 \(\mu\bigl(\partial B(x,r)\bigr)=0\)。

先补齐总质量可能无穷时所需的分解。
第\ref{b1:sec:0604}节的符号测度版本采用有限质量约定，不能不加说明就直接用于这里的 \(\nu\)。

\begin{proposition}[Lebesgue 分解，局部有限情形]\label{b1:prop:radon-lebesgue-decomposition}
存在非负 Borel 函数 \(f\in L^1_{\mathrm{loc}}(\mu)\) 及正 Radon 测度 \(\nu_s\)，使
\[
 \nu=f\mu+\nu_s,\quad \nu_s\perp\mu.
\]
其中 \(L^1_{\mathrm{loc}}(\mu)\) 指绝对值在每个紧集上可积。
密度 \(f\) 在 \(\mu\)-几乎处处相等的意义下唯一，\(\nu_s\) 唯一。
测度 \(\nu_a=f\mu\) 也是 Radon 测度。
\end{proposition}
\begin{proof}
令 \(\lambda=\mu+\nu\)。它是 \(\sigma\)-有限测度，故\cref{b1:thm:radon-nikodym}给出
\[
 \mu=p\lambda,\quad \nu=q\lambda,\quad p,q\geq0.
\]
因为 \(\lambda=(p+q)\lambda\)，密度唯一性给出 \(p+q=1\) 几乎处处。
修改一个 \(\lambda\)-零测 Borel 集上的值后，可使 \(p,q\) 处处非负且 \(p+q=1\)。
令
\[
 Z=\{\,x:p(x)=0\,\},\quad
 f(x)=
 \begin{cases}
 q(x)/p(x),&p(x)>0,\\
 0,&p(x)=0,
 \end{cases}
 \quad \nu_s(E)=\nu(E\cap Z).
\]
于是 \(\mu(Z)=\int_Zp\,\mathrm d\lambda=0\)，所以 \(\nu_s\perp\mu\)。
由\cref{b1:prop:density-change-integral}，对任意 Borel 集 \(E\) 有
\[
 \int_E f\,\mathrm d\mu
 =\int_E fp\,\mathrm d\lambda
 =\int_{E\setminus Z}q\,\mathrm d\lambda
 =\nu(E\setminus Z).
\]
这给出所需分解，且对紧集 \(K\)，有 \(\int_K f\,\mathrm d\mu\leq\nu(K)<\infty\)。

两个分量都不超过 \(\nu\)，因而局部有限。
若 \(\tau\leq\nu\)，对任意 Borel 集 \(E\) 及正整数 \(n\)，
在 \(E\cap\overline B(0,n)\) 内取紧集 \(K\)，使遗漏的 \(\nu\)-质量任意小，
则遗漏的 \(\tau\)-质量也任意小。
先逼近这一有界部分，再令 \(n\to\infty\)，即得 \(\tau\) 的紧内正则性。
应用于 \(\nu_a,\nu_s\)，便知二者都是 Radon 测度。

若还有 \(\nu=g\mu+\sigma\)，其中 \(\sigma\perp\mu\)，
取一个 \(\mu\)-零测 Borel 集 \(N\)，同时承载 \(\nu_s\) 与 \(\sigma\)。
在 \(\mathbb R^d\setminus N\) 上，\(f\mu=g\mu\)，密度唯一性给出 \(f=g\) 几乎处处。
而两个奇异分量在 \(N\) 上都等于 \(\nu\) 的限制，在 \(N\) 外都为零，故也相等。
这一论证没有对无穷总质量作减法。
\end{proof}

分解给出了一个积分意义的对象 \(f\)。本节的目标是证明
\[
 D_\mu\nu(x)=f(x)\quad\text{对 }\mu\text{-几乎处处的 }x.
\]
为此须分别处理 \(f\mu\) 的平均与 \(\nu_s\) 的比值；
前一项恢复密度，后一项在分母测度所见的几乎所有点消失。

\subsection{绝对连续部分}
\label{b1:subsec:140302}

下面采用 Fremlin 的《Measure Theory》\cite[472D]{Fremlin2016Measure}中
\term*[besicovitch-midu]{Besicovitch 密度定理}[Besicovitch density theorem]的证明路线。
它用细覆盖把局部平均过大的球合并成一项积分估计，适用于一般 Radon 测度。

\begin{theorem}[Besicovitch 密度]\label{b1:thm:radon-function-differentiation}
设 \(h\colon\mathbb R^d\to\mathbb K\) 为 Borel 函数，\(\mathbb K=\mathbb R\) 或 \(\mathbb C\)，
且 \(h\in L^1_{\mathrm{loc}}(\mu)\)。则对 \(\mu\)-几乎处处的 \(x\in S_\mu\)，有
\begin{equation}\label{b1:eq:radon-function-average}
 \lim_{r\downarrow0}
 \frac1{\mu\bigl(\overline B(x,r)\bigr)}
 \int_{\overline B(x,r)}h(y)\,\mathrm d\mu(y)=h(x),
\end{equation}
并且
\begin{equation}\label{b1:eq:radon-function-oscillation}
 \lim_{r\downarrow0}
 \frac1{\mu\bigl(\overline B(x,r)\bigr)}
 \int_{\overline B(x,r)}|h(y)-h(x)|\,\mathrm d\mu(y)=0.
\end{equation}
\end{theorem}
\begin{proof}
先设 \(h\) 为实值函数。记闭球平均为
\[
 A_rh(x)=\frac{\int_{\overline B(x,r)}h\,\mathrm d\mu}
                    {\mu\bigl(\overline B(x,r)\bigr)}
 \quad(x\in S_\mu).
\]
说明接下来所用坏集的可测性。
若 \(g\geq0\) 局部可积，则 \(g\mathbf1_{\overline B(0,n)}\mu\) 是有限 Radon 测度，
这由第\ref{b1:subsec:110204}小节的可积密度结论得到。
用这些限制穷竭可知 \(g\mu\) 局部有限且紧内正则，因而也是 Radon 测度。
分别取 \(g=h^+,h^-\)，\cref{b1:prop:relative-density-borel}的固定半径可测性便给出 \(A_rh\) 的 Borel 可测性。
在半径 \(r_j\downarrow r>0\) 时，闭球递减到 \(\overline B(x,r)\)，
局部可积性和有限测度下的连续性给出 \(A_{r_j}h(x)\to A_rh(x)\)。
因此在 \(0<r<1/k\) 内取平均的上确界或下确界时，可以只取有理半径。
相应的上、下极限都是 Borel 函数。

固定正整数 \(n\) 和有理数 \(a<b\)，令
\[
 E=\left\{\,x\in S_\mu\cap B(0,n):
 h(x)\leq a,\ \limsup_{r\downarrow0}A_rh(x)>b\,\right\}.
\]
这是有界 Borel 集，故 \(\mu(E)<\infty\)。给定 \(\varepsilon>0\)，
由\cref{b1:prop:integral-absolute-continuity}，可选 \(0<\eta<\varepsilon\)，使
\[
 F\subseteq B(0,n),\quad\mu(F)<\eta
 \quad\Longrightarrow\quad\int_F|h|\,\mathrm d\mu<\varepsilon.
\]
再由外正则性选开集 \(G\)，使
\[
 E\subseteq G\subseteq B(0,n),\quad\mu(G\setminus E)<\eta.
\]
对每个 \(x\in E\)，有任意小的中心闭球 \(B\subseteq G\)，满足
\(\int_Bh\,\mathrm d\mu>b\mu(B)\)。
将\cref{b1:cor:radon-vitali-covering}用于这些球，选得至多可数的两两不交闭球 \((B_j)\)，
使其并 \(H\) 满足 \(\mu(E\setminus H)=0\)。
因为 \(H\subseteq G\)，有 \(\mu(H\setminus E)<\eta\) 及
\(\int_{H\setminus E}|h|\,\mathrm d\mu<\varepsilon\)。
于是，不论 \(b\) 的正负，都有
\[
 \begin{aligned}
 b\mu(E)
 &\leq b\mu(H)+|b|\varepsilon
 =\sum_jb\mu(B_j)+|b|\varepsilon\\
 &\leq\sum_j\int_{B_j}h\,\mathrm d\mu+|b|\varepsilon
 =\int_Hh\,\mathrm d\mu+|b|\varepsilon\\
 &\leq\int_Eh\,\mathrm d\mu+(1+|b|)\varepsilon
 \leq a\mu(E)+(1+|b|)\varepsilon.
 \end{aligned}
\]
所有积分都在 \(B(0,n)\) 内，故有符号求和中没有无穷减无穷的问题。
令 \(\varepsilon\downarrow0\)，得到 \((b-a)\mu(E)=0\)，从而 \(\mu(E)=0\)。
对可数个 \(n,a,b\) 取并，便得
\[
 \limsup_{r\downarrow0}A_rh(x)\leq h(x)
 \quad\text{对 }\mu\text{-几乎处处的 }x.
\]
对 \(-h\) 应用同一结论，得到下极限不小于 \(h(x)\)，故\eqref{b1:eq:radon-function-average}成立。
复值情形分别对实部与虚部应用此结论，再取两个满测集的交即可。

最后证明绝对偏差的结论。
在 \(\mathbb R\) 情形取 \(Q=\mathbb Q\)，在 \(\mathbb C\) 情形取 \(Q=\mathbb Q+i\mathbb Q\)。
对每个 \(q\in Q\)，函数 \(g_q(y)=|h(y)-q|\) 局部可积。
由刚证得的平均公式及 \(Q\) 的可数性，有同一个满测集 \(D\subseteq S_\mu\)，使
\[
 \lim_{r\downarrow0}A_rg_q(x)=|h(x)-q|
 \quad(x\in D,\ q\in Q).
\]
固定 \(x\in D\) 及 \(\varepsilon>0\)，取 \(q\in Q\) 使 \(|h(x)-q|<\varepsilon\)。
当 \(r\) 充分小时，\(A_rg_q(x)<|h(x)-q|+\varepsilon\)，所以
\[
 \frac{\int_{\overline B(x,r)}|h(y)-h(x)|\,\mathrm d\mu(y)}
      {\mu\bigl(\overline B(x,r)\bigr)}
 \leq A_rg_q(x)+|q-h(x)|<3\varepsilon.
\]
令 \(\varepsilon\downarrow0\)，得到\eqref{b1:eq:radon-function-oscillation}。
\end{proof}

特别地，若 \(\nu\ll\mu\)，则上一小节的分解没有奇异分量，
\(\nu=f\mu\) 的局部质量比在 \(\mu\)-几乎处处收敛到 \(f\)。
密度只需局部可积，不需要连续、有界或相对于 Lebesgue 测度可积。
若改用 \(\mu\) 的完备化，可先取几乎处处相等的 Borel 代表；
零测集上修改函数既不改变球上的积分，也只改变零测集上的点值。

\subsection{奇异部分}
\label{b1:subsec:140303}

奇异测度未必只有点质量，承载它的零测集也可能稠密。
因此不能把证明简化成“几乎每一点的小球都避开奇异支集”。
所需的是球内奇异质量相对于分母质量趋于零，而不是最终完全没有奇异质量。

\begin{theorem}\label{b1:thm:radon-singular-differentiation}
若 \(\sigma\) 为 \(\mathbb R^d\) 上的正 Radon 测度且 \(\sigma\perp\mu\)，则
\[
 \lim_{r\downarrow0}
 \frac{\sigma\bigl(\overline B(x,r)\bigr)}
      {\mu\bigl(\overline B(x,r)\bigr)}=0
 \quad\text{对 }\mu\text{-几乎处处的 }x\in S_\mu.
\]
\end{theorem}
\begin{proof}
取 Borel 集 \(N\)，使 \(\mu(N)=0\)、\(\sigma(\mathbb R^d\setminus N)=0\)。
固定 \(t>0\) 和正整数 \(n\)，令
\[
 E=\{\,x\in(S_\mu\setminus N)\cap B(0,n):
              \overline D_\mu\sigma(x)>t\,\}.
\]
由\cref{b1:prop:relative-density-borel}，这是有界 Borel 集，而且 \(\sigma(E)=0\)。
给定 \(\varepsilon>0\)，外正则性给出开集 \(G\supseteq E\)，使 \(\sigma(G)<\varepsilon\)。
每个 \(x\in E\) 都有任意小的中心闭球 \(B\subseteq G\)，满足 \(\sigma(B)>t\mu(B)\)。
由\cref{b1:cor:radon-vitali-covering}，可选取这种球的可数不交族 \((B_j)\)，使其并在 \(\mu\)-意义下覆盖 \(E\)。
因此
\[
 t\mu(E)\leq t\sum_j\mu(B_j)
 \leq\sum_j\sigma(B_j)
 \leq\sigma(G)<\varepsilon.
\]
令 \(\varepsilon\downarrow0\)，得到 \(\mu(E)=0\)。
再对正整数 \(n\) 及正有理数 \(t\) 取可数并，便知上极限在 \(S_\mu\setminus N\) 上几乎处处为零。
比值非负，故极限也为零。
\end{proof}

若改用 \(\sigma\) 来选取“几乎处处”的点，结论的方向恰好相反。
由奇异性的对称性，把前一定理中的 \(\mu,\sigma\) 互换，得到
\[
 \lim_{r\downarrow0}
 \frac{\mu\bigl(\overline B(x,r)\bigr)}
      {\sigma\bigl(\overline B(x,r)\bigr)}=0
 \quad\text{对 }\sigma\text{-几乎处处的 }x\in S_\sigma.
\]
这里的分母始终为正。因而在这些点，把正数除以零规定为 \(+\infty\) 后，反向比值满足
\[
 \frac{\sigma\bigl(\overline B(x,r)\bigr)}
      {\mu\bigl(\overline B(x,r)\bigr)}\longrightarrow+\infty.
\]
具体地，对任意 \(M>0\)，当 \(r\) 足够小时，前一个比值小于 \(1/M\)；
若 \(\mu\) 的球质量为正，就取倒数，若它为零，反向比值已经是 \(+\infty\)。
这一扩展约定只用于 \(S_\sigma\) 上的反向结论，不为 \(0/0\) 赋值。

\subsection{Lebesgue 分解的微分解释}
\label{b1:subsec:140304}

现在把两项局部结论合并。

\begin{theorem}[Radon 测度的微分]\label{b1:thm:radon-measure-differentiation}
设 \(\mu,\nu\) 为 \(\mathbb R^d\) 上的正 Radon 测度，
\(\nu=\nu_a+\nu_s=f\mu+\nu_s\) 为\cref{b1:prop:radon-lebesgue-decomposition}中的分解。
则在 \(\mu\)-几乎处处的点，测度导数存在且有限，并满足
\[
 D_\mu\nu(x)
 =\lim_{r\downarrow0}
  \frac{\nu\bigl(\overline B(x,r)\bigr)}
       {\mu\bigl(\overline B(x,r)\bigr)}
 =\frac{\mathrm d\nu_a}{\mathrm d\mu}(x)=f(x).
\]
对 \(\nu_s\)-几乎处处的 \(x\)，同一比值在正数除以零取 \(+\infty\) 的约定下趋于 \(+\infty\)。
\end{theorem}
\begin{proof}
在 \(S_\mu\) 上，把球质量相加并除以正分母，得到
\[
 \frac{\nu\bigl(\overline B(x,r)\bigr)}
      {\mu\bigl(\overline B(x,r)\bigr)}
 =\frac{\int_{\overline B(x,r)}f\,\mathrm d\mu}
       {\mu\bigl(\overline B(x,r)\bigr)}
 +\frac{\nu_s\bigl(\overline B(x,r)\bigr)}
       {\mu\bigl(\overline B(x,r)\bigr)}.
\]
由\cref{b1:thm:radon-function-differentiation}，第一项几乎处处趋于 \(f(x)\)；
由\cref{b1:thm:radon-singular-differentiation}，第二项几乎处处趋于零。
取这两个满测集的交，再用 \(f\) 的局部可积性，即得第一个结论及极限有限性。
最后，上一小节已证 \(\nu_s\) 与 \(\mu\) 的反向比值在 \(\nu_s\)-几乎处处趋于无穷，
而 \(\nu\geq\nu_s\)，故 \(\nu\) 的相同比值也趋于无穷。
\end{proof}

例如在实线上取 \(\mu=m\)、\(\nu=m+\delta_0\)。当 \(x\ne0\) 时，足够小的闭区间不含原点，
故比值最终等于 \(1\)；在原点则有
\[
 \frac{\nu([-r,r])}{m([-r,r])}=1+\frac1{2r}\longrightarrow+\infty.
\]
于是 \(D_m\nu=1\) 在 \(m\)-几乎处处成立，但 \(\nu\) 并不绝对连续于 \(m\)。
单凭一个几乎处处有限的测度导数，不能断言奇异分量不存在。
正确的积分恢复公式是
\[
 \nu_a(E)=\int_E D_\mu\nu\,\mathrm d\mu,
\]
其中把导数未定义的 \(\mu\)-零测集上的值任意补成零。
右侧恢复的是绝对连续部分，不是自动恢复整个 \(\nu\)。

本节也说明了第\ref{b1:sec:1402}节结论中哪些内容依赖于 Lebesgue 测度，哪些不依赖。
有限维欧氏球的覆盖几何，足以保证任意 Radon 测度下的平均恢复与绝对偏差消失；
体积缩放律不是这一推广的前提。
至于在每个点选取怎样的函数代表，以及这些平均极限如何描述局部振荡，
下一节将在 Lebesgue 测度的背景下进一步讨论。

% !TeX root = ../../Book1.tex
\section{Lebesgue 点}
\label{b1:sec:1404}

% TODO: 撰写本节正文。

\subsection{Lebesgue 点定理}
\label{b1:subsec:140401}

待填充。

\subsection{精确代表}
\label{b1:subsec:140402}

待填充。

\subsection{局部平均}
\label{b1:subsec:140403}

待填充。

\subsection{\texorpdfstring{\(L^p\)}{Lp} 函数的 Lebesgue 点}
\label{b1:subsec:140404}

待填充。

% !TeX root = ../../Book1.tex
\section{近似极限与近似连续性}
\label{b1:sec:1405}

% TODO: 撰写本节正文。

\subsection{近似极限}
\label{b1:subsec:140501}

待填充。

\subsection{近似连续性}
\label{b1:subsec:140502}

待填充。

\subsection{密度拓扑的直观}
\label{b1:subsec:140503}

待填充。

\subsection{近似微分的预告}
\label{b1:subsec:140504}

待填充。

% 本章小结（不计入编号）
\begin{chaptersummary}
测度的微分把积分意义的密度转化为小球质量比，但这个转化有明确的取点基准：绝对连续密度在分母测度的几乎所有点恢复，奇异部分在那里消失；在奇异测度自己的几乎所有点，反向比例则显露出与基准测度不同的集中方式。因此，一个几乎处处有限的导数不等于没有奇异质量。

对于函数，Lebesgue 点控制平均误差的总量，精确代表从积分等价类选取点值，近似连续性则只控制大误差所占的比例。零测修改、跳跃与细小高峰分别说明这些概念不能混用。覆盖几何已经使一般欧氏 Radon 测度的微分成为可能；下一章将改变衡量集合大小的维数尺度，研究 Hausdorff 测度与维数，而不再只与环境体积比较。
\end{chaptersummary}


\BookExercises

% 八题均为自拟，以下注释只记录前置工具与题目目的，不冒称参考书原题。
% 正文来源对照：Fremlin2016Measure 472D（已实际读取）的一般Radon微分；
% Cantor测度使用本书0603已构造的数字映射，不预借第16章或未读书中的证明。

\begin{exercise}\label{b1:ex:14-one-sided-primitive}
设 \(f\in L^1((a,b))\)，并令
\[
 F(x)=\int_a^x f(t)\,\mathrm dt\quad(a\leq x\leq b).
\]
\begin{enumerate}
\item\label{b1:item:14-asymmetric-average}
若 \(x\in(a,b)\) 是 \(f\) 的 Lebesgue 点，
\(\alpha_r,\beta_r\geq0\)、\(\alpha_r+\beta_r>0\)，且 \(\alpha_r,\beta_r\to0\)，证明
\[
 \frac1{\alpha_r+\beta_r}
 \int_{x-\alpha_r}^{x+\beta_r}f(t)\,\mathrm dt\longrightarrow f(x).
\]
这里不要求 \(\alpha_r/\beta_r\) 有正的上下界，允许一端恒为零。
\item\label{b1:item:14-primitive-derivative}
证明 \(F\) 在 \([a,b]\) 上连续，并且在 \(f\) 的每个内部 Lebesgue 点都有 \(F'(x)=f(x)\)。
从而 \(F'=f\) 在 \((a,b)\) 上几乎处处成立。
\end{enumerate}
\end{exercise}
% 来源：自拟；把本章Lebesgue点的绝对偏差与第五章积分绝对连续性结合。
% 不预先使用“绝对连续函数必可微”或第16章微积分基本定理。
\begin{solution}
令 \(R_r=\max\{\alpha_r,\beta_r\}\)。所积分的区间包含在 \([x-R_r,x+R_r]\) 中，
且 \(R_r\leq\alpha_r+\beta_r\)。因此
\[
 \begin{aligned}
 \left|\frac1{\alpha_r+\beta_r}
       \int_{x-\alpha_r}^{x+\beta_r}f(t)\,\mathrm dt-f(x)\right|
 &\leq\frac1{\alpha_r+\beta_r}
       \int_{x-R_r}^{x+R_r}|f(t)-f(x)|\,\mathrm dt\\
 &\leq2\,\frac1{2R_r}
       \int_{x-R_r}^{x+R_r}|f(t)-f(x)|\,\mathrm dt
 \longrightarrow0.
 \end{aligned}
\]
这是一个一维特性：包含 \(x\) 的区间即使两端极不对称，长度仍能控制其到 \(x\) 的最远距离。

当 \(x<y\) 时，
\[
 |F(y)-F(x)|\leq\int_x^y|f(t)|\,\mathrm dt.
\]
由\cref{b1:prop:integral-absolute-continuity}，右侧随 \(y-x\to0\) 一致趋于零，
故 \(F\) 连续，包括两个端点。
若 \(x\) 是内部 Lebesgue 点，则 \(h>0\) 时的差商是 \([x,x+h]\) 上的平均，
\(h<0\) 时的差商是 \([x+h,x]\) 上的平均。
第一问分别给出两侧极限均为 \(f(x)\)，所以 \(F'(x)=f(x)\)。
最后应用\cref{b1:thm:lebesgue-point}即可得到几乎处处结论。
\end{solution}

\begin{exercise}\label{b1:ex:14-cantor-function}
令 \(\rho_C\) 为第\ref{b1:subsec:060303}小节用二进数字构造的 Cantor 概率测度，定义
\[
 F_C(x)=\rho_C\bigl((-\infty,x]\bigr)\quad(x\in\mathbb R).
\]
\begin{enumerate}
\item\label{b1:item:14-cantor-continuity}
证明 \(F_C\) 连续、单调不减，且 \(F_C(0)=0\)、\(F_C(1)=1\)。
进一步，令 \(\gamma=\log2/\log3\)，证明
\[
 |F_C(y)-F_C(x)|\leq4|y-x|^\gamma
 \quad(x,y\in\mathbb R).
\]
\item\label{b1:item:14-cantor-derivative}
证明 \(F_C'=0\) 在 Lebesgue 几乎处处成立，
但不存在 \(g\in L^1((0,1))\) 使 \(F_C(x)=\int_0^xg(t)\,\mathrm dt\) 对全部 \(x\in[0,1]\) 成立。
\end{enumerate}
函数 \(F_C|_{[0,1]}\) 称为\term[cantor-hanshu]{Cantor 函数}[Cantor function]。
\end{exercise}
% 来源：自拟；0603的数字映射给每个第n层Cantor基本区间质量2^{-n}。
% 从该已证测度构造分布函数，并与上一题的L1原函数比较；不调用BV理论。
\begin{solution}
单调性来自半直线的包含关系。
由第\ref{b1:subsec:060303}小节，\(\rho_C\) 无原子且全部质量在三分 Cantor 集 \(C\) 上。
测度从上、从下连续性分别给出
\[
 \lim_{y\downarrow x}F_C(y)=F_C(x),\quad
 \lim_{y\uparrow x}F_C(y)=\rho_C\bigl(( -\infty,x)\bigr)=F_C(x),
\]
所以 \(F_C\) 连续。端点值也由 \(\rho_C(C)=1\) 与 \(\rho_C(\{0\})=0\) 得到。

固定前 \(n\) 位三进数字，就确定一个长度 \(3^{-n}\) 的第 \(n\) 层基本闭区间。
在构造 \(\rho_C\) 的映射中，这等价于固定原像的前 \(n\) 位二进数字，
其原像为长度 \(2^{-n}\) 的二进半开区间，端点的编码差别不影响质量。
因此每个基本区间的 \(\rho_C\)-质量为 \(2^{-n}\)。
这些区间彼此之间的间隙至少为 \(3^{-n}\)，
故一个长度不超过 \(3^{-n}\) 的区间至多与两个基本区间有交。
若 \(0<y-x<1\)，取 \(n\geq0\)，使
\[
 3^{-(n+1)}<y-x\leq3^{-n}.
\]
于是
\[
 0\leq F_C(y)-F_C(x)
 =\rho_C\bigl((x,y]\bigr)
 \leq2\cdot2^{-n}
 <4|y-x|^\gamma.
\]
若 \(y-x\geq1\)，总质量不超过 \(1\) 已给出所需估计；交换 \(x,y\) 可处理另一种次序。

当 \(x\notin C\) 时，由 \(C\) 闭，可取一个不交于 \(C\) 的邻域。
测度在该邻域为零，\(F_C\) 因而在其中为常数，所以 \(F_C'(x)=0\)。
由于 \(m(C)=0\)，导数几乎处处为零。
假如存在题述 \(g\)，\cref{b1:ex:14-one-sided-primitive}给出 \(F_C'=g\) 几乎处处，
于是 \(g=0\) 几乎处处。这会推出 \(F_C(1)=0\)，与 \(F_C(1)=1\) 矛盾。
连续性、单调性以及几乎处处存在的导数，尚不足以保证由导数积分恢复原函数。
\end{solution}

\begin{exercise}\label{b1:ex:14-representative-null-change}
设 \(u\colon\mathbb R^d\to\mathbb R\) 连续，\(N\) 为 Lebesgue 零测集，
\(f=u\) 于 \(\mathbb R^d\setminus N\)，而在 \(N\) 上任意指定有限值。
\begin{enumerate}
\item\label{b1:item:14-representative-unique-value}
证明对每个 \(x\)，恰有一个数 \(a\) 满足
\[
 \lim_{r\downarrow0}\frac1{m\bigl(B(x,r)\bigr)}
       \int_{B(x,r)}|f(y)-a|\,\mathrm dy=0,
\]
且该数为 \(u(x)\)。据此求出 \(f\) 的 Lebesgue 点集及\cref{b1:def:precise-representative}中的精确代表。
\item\label{b1:item:14-representative-unbounded}
构造一个有限值 Borel 函数 \(f\)，使它在每个非空开集中都无界，
但其精确代表处处为零，且几乎每一点都是 \(f\) 的 Lebesgue 点。
\end{enumerate}
\end{exercise}
% 来源：自拟；区分平均定义的唯一值、指定代表的点值与通常连续性。
% 第二问把普通的可数集改值例增强为处处局部无界，避免只复述正文的零测修改规则。
\begin{solution}
由于 Lebesgue 测度完备，\(f\) 可测且与 \(u\) 具有相同的局部积分。
连续性给出
\[
 \frac1{m\bigl(B(x,r)\bigr)}\int_{B(x,r)}|f(y)-u(x)|\,\mathrm dy
 =\frac1{m\bigl(B(x,r)\bigr)}\int_{B(x,r)}|u(y)-u(x)|\,\mathrm dy
 \longrightarrow0.
\]
若 \(a\) 也满足题述极限，三角不等式给出
\[
 |a-u(x)|
 \leq\frac1{m\bigl(B(x,r)\bigr)}\int_{B(x,r)}|a-f(y)|\,\mathrm dy
 +\frac1{m\bigl(B(x,r)\bigr)}\int_{B(x,r)}|f(y)-u(x)|\,\mathrm dy.
\]
令 \(r\downarrow0\)，得 \(a=u(x)\)。
因此 \(x\) 是指定代表 \(f\) 的 Lebesgue 点，当且仅当 \(f(x)=u(x)\)。
另一方面，\(f\) 的球平均在每一点都趋于 \(u(x)\)，故其精确代表处处等于 \(u\)。

第二问中，设 \(e_1\) 为第一坐标向量，并取
\[
 D_k=\mathbb Q^d+k\sqrt2\,e_1\quad(k\geq1).
\]
每个 \(D_k\) 可数且稠密。若 \(D_k\cap D_l\ne\varnothing\)，
第一坐标会给出 \((k-l)\sqrt2\in\mathbb Q\)，所以 \(k=l\)；这些集合两两不交。
令 \(f=k\) 于 \(D_k\)，其余点取 \(0\)。
它是有限值 Borel 函数，非零点集 \(D=\bigcup_kD_k\) 可数，因而 \(f=0\) 几乎处处。
任意非空开集都与每个 \(D_k\) 相交，所以 \(f\) 在其中无界，特别地处处不连续。
将第一问用于 \(u=0\)，可知其精确代表处处为零，Lebesgue 点集恰为 \(\mathbb R^d\setminus D\)。
\end{solution}

\begin{exercise}\label{b1:ex:14-products-at-lebesgue-points}
设 \(1<p,q<\infty\)、\(1/p+1/q=1\)，
\(f\in L^p_{\mathrm{loc}}(\mathbb R^d)\)、\(g\in L^q_{\mathrm{loc}}(\mathbb R^d)\)。
\begin{enumerate}
\item\label{b1:item:14-product-holder-point}
若在一点 \(x\) 有
\[
 \begin{aligned}
 \frac1{m\bigl(B(x,r)\bigr)}\int_{B(x,r)}|f(y)-f(x)|^p\,\mathrm dy&\longrightarrow0,\\
 \frac1{m\bigl(B(x,r)\bigr)}\int_{B(x,r)}|g(y)-g(x)|^q\,\mathrm dy&\longrightarrow0,
 \end{aligned}
\]
证明 \(x\) 是乘积 \(fg\) 的 Lebesgue 点。
\item\label{b1:item:14-product-failure}
构造实线上非负函数 \(h\in L^2(\mathbb R)\)，使 \(h(0)=0\)、\(0\) 是 \(h\) 的 Lebesgue 点，
但不是 \(h^2\) 的 Lebesgue 点。由此说明“两个函数都有同一个 Lebesgue 点”不足以保证乘积也有该 Lebesgue 点，
即使乘积已经局部可积。
\end{enumerate}
\end{exercise}
% 来源：自拟；第一问用第10章Hölder/Minkowski，第二问采用显式尖峰。
% 反例同时区分L1平均偏差与L2平均偏差，不误把逐点的L1性质提升为Lp性质。
\begin{solution}
对球 \(B=B(x,r)\)，把 \(m(B)^{-1}m|_B\) 视为概率测度。
由 \(fg-f(x)g(x)=(f-f(x))g+f(x)(g-g(x))\) 及 Hölder 不等式，
\[
 \begin{aligned}
 \frac1{m(B)}\int_B|fg-f(x)g(x)|
 &\leq
 \left(\frac1{m(B)}\int_B|f-f(x)|^p\right)^{1/p}
 \left(\frac1{m(B)}\int_B|g|^q\right)^{1/q}\\
 &\quad+|f(x)|\left(\frac1{m(B)}\int_B|g-g(x)|^q\right)^{1/q}.
 \end{aligned}
\]
Minkowski 不等式还给出
\[
 \left(\frac1{m(B)}\int_B|g|^q\right)^{1/q}
 \leq|g(x)|+
 \left(\frac1{m(B)}\int_B|g-g(x)|^q\right)^{1/q},
\]
故这一因子保持有界，前一式右侧趋于零。局部可积性 \(fg\in L^1_{\mathrm{loc}}\) 也由 Hölder 不等式得到。

第二问取两两不交区间及函数
\[
 I_n=(2^{-n},2^{-n}+2^{-3n}),\quad
 h=\sum_{n=2}^{\infty}2^n\mathbf1_{I_n}.
\]
于是 \(h(0)=0\)，且
\[
 \int_{\mathbb R}h^2=\sum_{n=2}^{\infty}2^{2n}2^{-3n}
 =\sum_{n=2}^{\infty}2^{-n}<\infty.
\]
当 \(2^{-(N+1)}<r\leq2^{-N}\) 时，\((-r,r)\) 只能与指标 \(n\geq N+1\) 的区间有正测度交集。
因此
\[
 \frac1{2r}\int_{-r}^rh
 \leq\frac1{2r}\sum_{n=N+1}^{\infty}2^{-2n}
 \leq\frac13\,2^{-N}\longrightarrow0.
\]
故 \(0\) 是 \(h\) 的 Lebesgue 点。
但取 \(r_n=2^{-n}+2^{-3n}\)，整个 \(I_n\) 都落在 \((-r_n,r_n)\) 内，从而
\[
 \frac1{2r_n}\int_{-r_n}^{r_n}|h(t)^2-h(0)^2|\,\mathrm dt
 \geq\frac{2^{-n}}{2(2^{-n}+2^{-3n})}
 \longrightarrow\frac12.
\]
所以 \(0\) 不是 \(h^2\) 的 Lebesgue 点。
\end{solution}

\begin{exercise}\label{b1:ex:14-finite-p-endpoint}
令
\[
 J_n=(2^{-n},2^{-n}+2^{-2n-2})\quad(n\geq2),\quad
 E=\bigcup_{n=2}^{\infty}J_n,\quad f=\mathbf1_E.
\]
证明对每个有限 \(p\geq1\)，都有
\[
 \lim_{r\downarrow0}\frac1{2r}\int_{-r}^r|f(t)-f(0)|^p\,\mathrm dt=0,
\]
但是对每个 \(r>0\)，都有
\[
 \operatorname*{ess\,sup}_{|t|<r}|f(t)-f(0)|=1.
\]
进一步说明，即使 \(f\in L^\infty\)，全部有限幂次的平均偏差结论也不推出本质一致意义的局部收敛。
\end{exercise}
% 来源：自拟；有界示性函数的稀薄区间例检验p=∞端点，不重复正文的有限p证明。
\begin{solution}
由于 \(0\notin E\)，\(f(0)=0\)，且 \(|f|^p=\mathbf1_E\) 对每个 \(p>0\) 都成立。
若 \(2^{-(N+1)}<r\leq2^{-N}\)，则只有 \(n\geq N+1\) 的区间可能在 \((-r,r)\) 内贡献正长度，故
\[
 m\bigl(E\cap(-r,r)\bigr)
 \leq\sum_{n=N+1}^{\infty}2^{-2n-2}
 =\frac1{12}\,2^{-2N}.
\]
除以 \(2r>2^{-N}\)，得到平均值不超过 \(2^{-N}/12\)，趋于零。
这个估计与有限指数 \(p\) 无关。

另一方面，对任意 \(r>0\)，可取充分大的 \(n\)，使整个 \(J_n\subset(-r,r)\)。
在这个正长度区间上 \(f=1\)，所以本质上确界至少为 \(1\)；又因 \(0\leq f\leq1\)，它恰为 \(1\)。
有限幂次的平均允许高度为 \(1\) 的偏差占据越来越小的相对体积，而本质上确界不会忽略正测度的尖峰。
\end{solution}

\begin{exercise}\label{b1:ex:14-affine-lebesgue-points}
设 \(T(x)=Ax+b\)，其中 \(A\) 为可逆实 \(d\times d\) 矩阵，\(b\in\mathbb R^d\)，
并设 \(f\in L^1_{\mathrm{loc}}(\mathbb R^d)\)。
\begin{enumerate}
\item\label{b1:item:14-affine-invariance}
证明 \(f\circ T\) 局部可积，并且 \(x\) 是 \(f\circ T\) 的 Lebesgue 点，当且仅当 \(T(x)\) 是 \(f\) 的 Lebesgue 点。
特别地，说明平移如何改变 Lebesgue 点集。
\item\label{b1:item:14-singular-pullback}
设 \(L\colon\mathbb R\to\mathbb R^2\)、\(L(t)=(t,0)\)。
构造两个在二维 Lebesgue 几乎处处相等的 Borel 函数 \(f,g\)，
使 \(f\circ L\) 与 \(g\circ L\) 在每一点都不相等。
解释为何非可逆映射的复合不能无条件定义在 Lebesgue 几乎处处等价类上。
\end{enumerate}
\end{exercise}
% 来源：自拟；第7章线性换元与本章局部平均结合；奇异线性映射一问检查零集逆像。
\begin{solution}
若 \(K\subset\mathbb R^d\) 紧，线性换元给出
\[
 \int_K|f\bigl(T(z)\bigr)|\,\mathrm dz
 =\frac1{|\det A|}\int_{T(K)}|f(y)|\,\mathrm dy<\infty,
\]
因为 \(T(K)\) 紧，故 \(f\circ T\) 局部可积。
令 \(y_0=T(x)\)，并取 \(c=\|A\|>0\)，使 \(|Av|\leq c|v|\)。
由 \(T(B(x,r))\subseteq B(y_0,cr)\)，有
\[
 \begin{aligned}
 &\frac1{m\bigl(B(x,r)\bigr)}
    \int_{B(x,r)}|f\bigl(T(z)\bigr)-f(y_0)|\,\mathrm dz\\
 &\quad\leq
 \frac1{|\det A|\,m\bigl(B(x,r)\bigr)}
    \int_{B(y_0,cr)}|f(y)-f(y_0)|\,\mathrm dy\\
 &\quad=
 \frac{c^d}{|\det A|}\,
 \frac1{m\bigl(B(y_0,cr)\bigr)}
    \int_{B(y_0,cr)}|f(y)-f(y_0)|\,\mathrm dy.
 \end{aligned}
\]
若 \(y_0\) 是 \(f\) 的 Lebesgue 点，右侧趋于零，便得一个方向。
对 \(T^{-1}\) 与函数 \(f\circ T\) 使用相同论证，得到反向。
所以 Lebesgue 点集按 \(T^{-1}\) 变换；当 \(T(x)=x+b\) 时，它平移为原点集减去 \(b\)。

第二问取 \(f=\mathbf1_{\mathbb R\times\{0\}}\)、\(g=0\)。
横轴是二维零测集，所以 \(f=g\) 几乎处处。
但 \(f\circ L=1\)、\(g\circ L=0\) 在整条实线上成立。
这里一个二维零测集的逆像是整个定义域，而不再是一维零测集。
因此在没有相应零集逆像条件时，复合的结果取决于所选代表。
\end{solution}

\begin{exercise}\label{b1:ex:14-degenerate-rectangles}
在 \(\mathbb R^2\) 中令
\[
 E=\{\,(x,y):|x|<1,\ |y|<x^2\,\},\quad f=\mathbf1_E,
 \quad R_r=[-r,r]\times[-r^2,r^2]\quad(0<r<1).
\]
证明原点是 \(f\) 的 Lebesgue 点且 \(f(0,0)=0\)，但是
\[
 \frac1{m_2(R_r)}\int_{R_r}f\,\mathrm dm_2=\frac13
 \quad(0<r<1).
\]
结合\cref{b1:prop:lebesgue-controlled-sets}，指出这些集合虽然以原点为中心且直径趋于零，却缺少了哪一项控制。
\end{exercise}
% 来源：自拟；尖区面积由第7章Tonelli直接计算；该反例留在习题，不复述正文示例。
\begin{solution}
由定义，原点不在 \(E\)。当 \(0<r<1\) 时，球内的横坐标满足 \(|x|<r\)，故
\[
 m_2\bigl(E\cap B(0,r)\bigr)
 \leq\int_{-r}^r2x^2\,\mathrm dx=\frac43r^3.
\]
球面积为 \(\pi r^2\)，所以 \(f\) 在球内的绝对偏差平均不超过 \(4r/(3\pi)\)，趋于零。
于是原点是值为零的 Lebesgue 点。

在 \(R_r\) 内，对每个 \(|x|<r\) 都有 \(x^2\leq r^2\)，
因而 \(E\) 的纵向截面仍有完整长度 \(2x^2\)。所以
\[
 \int_{R_r}f\,\mathrm dm_2=\int_{-r}^r2x^2\,\mathrm dx=\frac43r^3,
 \quad m_2(R_r)=4r^3,
\]
相除得到 \(1/3\)。
矩形的长短边之比为 \(1/r\)，趋于无穷；包含它的中心球半径可取 \(\sqrt{r^2+r^4}\)。
对应的体积比为
\[
 \frac{m_2(R_r)}{m_2\bigl(B(0,\sqrt{r^2+r^4})\bigr)}
 =\frac{4r}{\pi(1+r^2)}\longrightarrow0.
\]
因此不存在一个正的常数，统一从下控制矩形体积与外接球体积之比。
细长矩形把越来越大比例的自身面积放在球内的稀薄部分上，球平均结论便不能传递过去。
\end{solution}

\begin{exercise}\label{b1:ex:14-accumulating-atoms}
在实线上定义
\[
 \mu=\sum_{n=1}^{\infty}2^{-n}\delta_{2^{-n}},\quad
 \nu=\sum_{k=1}^{\infty}2^{-2k}\delta_{2^{-2k}}.
\]
\begin{enumerate}
\item\label{b1:item:14-atoms-radon}
证明 \(\mu,\nu\) 都是有限 Radon 测度，\(\nu\ll\mu\)，并给出一个 Borel 密度 \(h\)，使 \(\nu=h\mu\)。
求出 \(\operatorname{supp}\mu\)。
\item\label{b1:item:14-atoms-derivative}
在每个原子 \(2^{-n}\) 处计算 \(D_\mu\nu\)。再证明，虽然 \(0\in\operatorname{supp}\mu\)，
但 \(D_\mu\nu(0)\) 不存在。
解释这为何不与\cref{b1:thm:radon-measure-differentiation}矛盾。
\end{enumerate}
\end{exercise}
% 来源：自拟；几何级数精确计算一个支集内的例外点，区分“支集上处处”和“测度几乎处处”。
\begin{solution}
两者总质量分别为 \(1\) 与 \(1/3\)，且 \(0\leq\nu\leq\mu\)。
它们的可数可加性来自非负级数换序。
对任意 Borel 集，只保留前 \(N\) 个原子，得到包含于该集的有限点集；
遗漏质量至多为 \(\sum_{n>N}2^{-n}\)，趋于零。
有限点集紧，故这两个测度紧内正则，且总质量有限保证局部有限，所以都是 Radon 测度。
取
\[
 h=\mathbf1_{\{\,2^{-2k}:k\geq1\,\}},
\]
即可逐项验证 \(\nu=h\mu\)，从而 \(\nu\ll\mu\)。
每个原子都在支集中；原点的任意邻域也包含尾部原子并具有正质量。
而集合 \(S=\{0\}\cup\{\,2^{-n}:n\geq1\,\}\) 闭，其外每一点都有不含原子的邻域。
因此 \(\operatorname{supp}\mu=S\)。

固定 \(n\)，当 \(r\) 小于该原子与其他原子之间距离的一半时，
\(\overline B(2^{-n},r)\) 只含这个原子。
所以球比值最终恒等于 \(1\)（\(n\) 偶数）或 \(0\)（\(n\) 奇数），恰为 \(h(2^{-n})\)。

在原点取 \(r_N=2^{-N}\)。此时
\[
 \mu([-r_N,r_N])=\sum_{n=N}^{\infty}2^{-n}=2^{1-N},
\]
而分子的求和只保留偶数指标，因而
\[
 \frac{\nu([-r_N,r_N])}{\mu([-r_N,r_N])}
 =\begin{cases}
 2/3,&N\text{ 偶数},\\
 1/3,&N\text{ 奇数}.
 \end{cases}
\]
两个子列的极限不同，故 \(D_\mu\nu(0)\) 不存在。
但 \(\mu(\{0\})=0\)，而全部正质量原子处的公式都成立，故这只是定理所允许的零测例外。
原点属于支集，只保证每个球的分母为正，并不保证该点的微分极限存在。
\end{solution}
